\sloppy
\chapter{Introduction}

Fake news detection is the automated task of identifying false, misleading, or fabricated information that is presented as genuine news, whether the deception is deliberate or not. It relies on methods from Natural Language Processing (NLP) and Machine Learning (ML)~\citep{Shu2020FakeNewsNet,Shu2017FakeNewsDetection}. Digital media platforms have sharply increased the volume, speed, and reach of the information shared online every day. News, opinion, and rumour travel through the same channels and often look alike at first glance, so readers must judge credibility from a headline and a few seconds of attention.

The consequences are serious. Fake news can manipulate public perception, interfere with democratic processes such as elections, and gradually erode confidence in trusted sources~\citep{Shu2020FakeNewsNet,Shu2017FakeNewsDetection}. Empirical studies show that false information spreads faster and reaches a wider audience on social media than verified news~\citep{Vosoughi2018Spread}. At that pace, journalists and professional fact-checkers cannot verify every claim before it gets out of control.

The harm also goes beyond individual misunderstanding. Large-scale disinformation campaigns have been linked to social unrest, financial market manipulation, and public health harms. Building systems that detect fake news reliably, at scale, and in near real time is therefore an urgent research problem, and it motivates the work in this thesis.

\section{Challenges and Issues}

``Fake news'' is not one homogeneous kind of content. It is an umbrella term for several classes of information, each with its own motive, structure, and linguistic signature. The literature commonly identifies five recurring types:

\begin{itemize}
\item \textbf{False News Stories:} Entirely fabricated information designed to deceive readers.
\item \textbf{Manipulated News:} Genuine information that has been selectively edited, distorted, or taken out of context.
\item \textbf{Satirical News:} Exaggerated or humorous content that is not meant to deceive but may be mistaken for factual reporting.
\item \textbf{Clickbait:} Stories with sensational or misleading headlines that do not match their content.
\item \textbf{Propaganda:} Biased or slanted information produced to advance a particular agenda or cause.
\end{itemize}

This diversity makes automatic detection hard. A model trained on the markers of one category, such as sensational clickbait headlines, may fail on subtly manipulated but factually grounded reporting, or on satire that imitates a legitimate outlet's tone. Detection therefore needs models that capture deep contextual and semantic relationships in text, and not only superficial keyword or stylistic cues.

The modelling side adds a second difficulty. Classical machine learning algorithms such as Logistic Regression, Support Vector Machines (SVM), Naive Bayes, and Random Forest have long been applied to this task with hand-engineered statistical features such as TF-IDF vectors. They are simple, interpretable, and cheap to run, but they cannot capture word order or long-range semantic dependency, and this limits their accuracy. Sequential deep learning models such as Long Short-Term Memory (LSTM) and Bidirectional LSTM (BiLSTM) networks were introduced to fill that gap by modelling word sequences directly. Their effectiveness can still suffer when contextual relationships are complex, training data is limited, or training time and hyperparameter tuning are restricted. The experiments in this thesis show exactly this limitation (Chapter~4).

Transformer-based models, most notably BERT (Bidirectional Encoder Representations from Transformers), outperform these earlier sequential approaches by producing deeply bidirectional, contextualized representations of text in a single forward pass~\citep{Kaliyar2021FakeBERT,Mehta2022RecentTrends}. Sequential models read text step by step, whereas BERT attends to the whole input at once through self-attention and so sees the relationships between words more fully. The contrast between sequential and attention-based modelling is one of the central technical questions of this thesis.

A last challenge concerns evaluation. Benchmark datasets such as ISOT carry source artifacts, for example a news agency's name in the headline, and a model can score highly by recognising these markers instead of the content. A model that does well on one dataset may therefore fail on another, so results from a single dataset can overstate real performance. This thesis examines the question with a cross-dataset test.

\section{Motivation of the Research}

The motivation is both societal and technical. Societally, misinformation spreads at a scale and speed that manual fact-checking alone cannot match. An automated, scalable first line of detection is needed to support human fact-checkers and platform moderators, not to replace them. Technically, the field has moved quickly from classical statistical text representations (Bag-of-Words, TF-IDF), through static embeddings (Word2Vec, GloVe), to deep contextual embeddings (BERT). For a given dataset and task, it is often unclear whether the added complexity of transformer and hybrid architectures buys a matching gain in accuracy.

This thesis therefore compares a purely sequential deep learning model (BiLSTM), a purely transformer-based model (BERT), and a hybrid of the two (BERT + BiLSTM) across several benchmark datasets. The comparison makes the trade-offs between model complexity, training cost, and detection performance visible, and it informs deployment decisions for real-world systems such as social media monitoring pipelines and automated fact-checking assistants.

\section{Research Objectives}

The objectives of this research are:

\begin{enumerate}
\item To design a multi-stage fake news detection pipeline that combines NLP-based text preprocessing, several feature extraction and representation techniques, and a classification layer.
\item To implement and train three classification models, a standalone BiLSTM network, a fine-tuned BERT model, and a hybrid BERT + BiLSTM model, under comparable training conditions.
\item To evaluate and compare the three models with standard classification metrics (accuracy, precision, recall, and F1-score) and with visual diagnostics (confusion matrices, ROC curves, precision--recall curves, and loss curves).
\item To assess how well the models generalize across the benchmark datasets LIAR and ISOT, including a cross-dataset test in which a model trained on one dataset is evaluated on the other.
\item To identify the practical and architectural trade-offs between sequential, transformer-based, and hybrid models for fake news detection, and to outline directions for future improvement using large language models and explainable AI.
\end{enumerate}

\section{Thesis Contributions}

This thesis makes the following contributions:

\begin{itemize}
\item It proposes a layered architecture for fake news detection that combines NLP preprocessing, several feature extraction techniques, and a comparative classification framework.
\item It presents a systematic experimental comparison of BiLSTM, BERT, and a hybrid BERT + BiLSTM model across several benchmark datasets, with GRU-based hybrids and a zero-shot local LLM as additional baselines.
\item It shows, with quantitative evidence, that transformer-based models outperform purely sequential deep learning models on this task under constrained training budgets.
\item It documents the limitations of sequential models and their hybrid variants when trained for few epochs and without extensive hyperparameter tuning.
\item It reports a cross-dataset evaluation on cleaned, deduplicated headlines, which shows that models scoring above 95\% on ISOT fall to about chance level on LIAR2.
\item It identifies promising directions for future research, including Large Language Models (LLMs), Retrieval-Augmented Generation (RAG), and Explainable AI (XAI) in fake news detection pipelines.
\end{itemize}

\section{Thesis Organization}

The remaining chapters are organized as follows.

\subsubsection*{Chapter 1: Introduction}

This chapter introduced the problem of fake news detection and its societal significance, the challenges of separating genuine from fabricated content automatically, the motivation for this research, and the objectives and contributions of the thesis.

\subsubsection*{Chapter 2: Literature Survey}

This chapter reviews the existing literature on fake news detection. It runs from classical machine learning techniques, through deep learning and hybrid architectures, to transformer-based and evidence-aware methods, benchmark dataset papers, and recent work on large language models and explainability.

\subsubsection*{Chapter 3: Proposed Approach}

This chapter describes the proposed multi-stage detection system: the problem statement, the overall system model and pipeline, the BERT + BiLSTM architecture, the underlying algorithm, and the reasons for combining a transformer encoder with a sequential BiLSTM layer.

\subsubsection*{Chapter 4: Experimental Results}

This chapter presents the experimental setup, datasets, evaluation metrics, and hyperparameter configuration used for the three models. It then reports and discusses the comparative results, confusion matrices, ROC curves, precision--recall curves, and training and validation loss curves, followed by the GRU experiments, the zero-shot LLM test, and the cross-dataset evaluation.

\subsubsection*{Chapter 5: Conclusion and Future Direction}

This chapter summarizes the main findings, discusses what they mean for practical fake news detection systems, and outlines future work on large language models, Retrieval-Augmented Generation, Explainable AI, multilingual and multi-modal detection, and federated deployment.
